% ============================================================
%  Beating Shannon: encryption from a PRG
% ============================================================

\section{Beating Shannon: encryption from a PRG}\label{sec:prg-encryption}

We now have all the tools to go back to the problem that opened this part of the course. Shannon's theorem (\Cref{thm:shannon}) says that a perfectly secret scheme needs a key at least as long as the message. The One-Time Pad meets this bound, but a gigabyte of message requires a gigabyte of key. With a PRG we can \textbf{stretch} a short key into a long pad, and give up only perfect secrecy against unbounded adversaries, keeping security against efficient ones.

\subsection{The scheme}

\begin{frameddefn}[label={def:prg-otp}]{Encryption from a PRG}
  Let $G : \{0,1\}^n \to \{0,1\}^{2n}$ be a PRG. The scheme $\Pi = (\Enc, \Dec)$ has key space $\Kspace = \{0,1\}^n$ and message and ciphertext spaces $\Mspace = \Cspace = \{0,1\}^{2n}$:
  \[
    \Enc(k, m) = G(k) \oplus m \qquad\qquad \Dec(k, c) = G(k) \oplus c
  \]
  with the key $k \sample \{0,1\}^n$ chosen uniformly.
\end{frameddefn}

\begin{figure}[H]
  \centering
  \begin{tikzpicture}[box/.style={draw, rounded corners=2pt, minimum height=7mm, inner sep=4pt}]
    \node (k) at (0,0) {$k$};
    \node[box] (G) at (2,0) {$G$};
    \node[draw, circle, inner sep=1pt] (x) at (5,0) {$\oplus$};
    \node (m) at (5,1.3) {$m$};
    \node (c) at (7.5,0) {$c$};
    \draw[msg] (k) -- node[above, lbl] {$n$ bits} (G);
    \draw[msg] (G) -- node[above, lbl] {$G(k)$, $2n$ bits} (x);
    \draw[msg] (m) -- (x);
    \draw[msg] (x) -- (c);
  \end{tikzpicture}
  \caption{It is the One-Time Pad, with the pad $G(k)$ produced by the PRG instead of being a truly random key.}
\end{figure}

Observations:
\begin{itemize}
  \item \textbf{Correctness} $\to$ $\Dec(k, \Enc(k, m)) = G(k) \oplus G(k) \oplus m = m$.
  \item \textbf{The key is shorter than the message} $\to$ $n$ bits of key for $2n$ bits of message. With the stretching construction (\Cref{thm:stretch}) the message can be any polynomial length $n + \ell$.
  \item \textbf{No contradiction with Shannon} $\to$ the scheme is \textbf{not} perfectly secret: an unbounded adversary can try all the $2^n$ keys. The security we will prove holds only against efficient adversaries.
\end{itemize}

\begin{framedobs}{It is still a one-time scheme}
  As with the One-Time Pad, the key can be used \textbf{only once}. If two messages are encrypted with the same key,
  \[
    c_1 = G(k) \oplus m_1, \qquad c_2 = G(k) \oplus m_2 \qquad \Longrightarrow \qquad c_1 \oplus c_2 = m_1 \oplus m_2
  \]
  The pad cancels out, and Eve learns the XOR of the two messages without knowing the key. If she knows (or guesses) $m_1$, she recovers $m_2 = c_1 \oplus c_2 \oplus m_1$; and since real messages are very redundant (natural language, file headers), $m_1 \oplus m_2$ alone is usually enough to recover both.
\end{framedobs}

\subsection{Defining computational security}

To state what the scheme achieves we need a \textbf{computational} version of perfect secrecy. Recall the formulation of perfect secrecy in terms of distributions (statement (iii) of \Cref{lem:equivalent}): for every pair of messages $m_0, m_1$,
\[
  \Enc(K, m_0) \equiv \Enc(K, m_1) \qquad \text{with } K \text{ uniform}
\]
that is, the two ciphertexts have \textbf{exactly the same distribution}, so not even an unbounded adversary can tell them apart. The computational version replaces ``identical'' ($\equiv$) with ``computationally indistinguishable'' ($\cind$). It is convenient to write it as a game with the bit $b$ \textbf{fixed}, and to compare the two versions of the game.

\begin{frameddefn}[label={def:one-time-game}]{The one-time game $\mathrm{GAME}^{\text{one-time}}_{\Pi,\Adv}(\lambda, b)$}
  Let $\lambda$ be the security parameter (the length $n$ of the key) and $b \in \{0,1\}$ a fixed bit. The adversary $\Adv$ plays against a challenger $C$:
  \begin{enumerate}
    \item $\Adv$ chooses two messages $m_0, m_1 \in \Mspace$ of the same length and sends them to $C$;
    \item $C$ chooses a key $k \sample \Kspace$ and sends back $c = \Enc(k, m_b)$;
    \item $\Adv$ outputs a bit $b'$, which is the output of the game.
  \end{enumerate}
\end{frameddefn}

\begin{center}
  \begin{tikzpicture}[x=1cm, y=1cm]
    \node (A) at (0,0) {$\Adv$};
    \node (C) at (6,0) {$C$};
    \draw[msg] (0,-0.6) -- node[above] {$m_0, m_1 \in \Mspace$} (6,-0.6);
    \node[lbl, anchor=west] at (6.6,-1.0) {$k \sample \Kspace$};
    \node[lbl, anchor=west] at (6.6,-1.45) {$c = \Enc(k, m_b)$};
    \draw[msg] (6,-1.4) -- node[above] {$c$} (0,-1.4);
    \draw[msg] (0,-2.2) -- node[above] {$b'$} (6,-2.2);
  \end{tikzpicture}
\end{center}

\begin{frameddefn}[label={def:one-time-security}]{One-time computational security}
  A symmetric-key encryption scheme $\Pi = (\Enc, \Dec)$ is \textbf{one-time computationally secure} if
  \[
    \mathrm{GAME}^{\text{one-time}}_{\Pi,\Adv}(\lambda, 0) \cind \mathrm{GAME}^{\text{one-time}}_{\Pi,\Adv}(\lambda, 1)
  \]
  that is, for every PPT adversary $\Adv$ there is a negligible function $\negl$ such that
  \[
    \Big|\Pr\big[\mathrm{GAME}^{\text{one-time}}_{\Pi,\Adv}(\lambda, 0) = 1\big] - \Pr\big[\mathrm{GAME}^{\text{one-time}}_{\Pi,\Adv}(\lambda, 1) = 1\big]\Big| \leq \negl(\lambda)
  \]
\end{frameddefn}

In words: whatever pair of messages Eve chooses, she cannot tell whether she is looking at an encryption of $m_0$ or of $m_1$, except with negligible advantage.

\paragraph{The same definition as in Chapter 1.} This is equivalent to the definition of \Cref{def:eav}, where the bit $b$ is chosen at random and the adversary must guess it. Writing $p_b = \Pr[\mathrm{GAME}(\lambda, b) = 1]$ for the probability that $\Adv$ outputs $1$ when the bit is $b$:
\[
  \Pr[b' = b] = \frac{1}{2}\Pr[b' = 0 \mid b = 0] + \frac{1}{2}\Pr[b' = 1 \mid b = 1] = \frac{1}{2}(1 - p_0) + \frac{1}{2}p_1 = \frac{1}{2} + \frac{1}{2}(p_1 - p_0)
\]
So guessing $b$ with probability $\frac{1}{2} + \negl$ is the same as having $|p_1 - p_0| \leq \negl$, up to the factor $2$ (and to flipping the answer, which turns a negative difference into a positive one).

\subsection{Sanity checks}

A good definition must imply everything we intuitively expect from a secure scheme. Let's check.

\begin{itemize}
  \item \textbf{The key is hard to compute.} If an adversary could compute $k$ from $c$, it would win the game: choose two different messages $m_0, m_1$, compute $k$, decrypt $c$ and output $b' = 0$ if the result is $m_0$, and $1$ otherwise. Its advantage would be $1$.

  \textbf{But hiding the key is not enough.} The scheme $\Enc(k, m) = m$ hides the key perfectly (the ciphertext does not even depend on it), yet it reveals the whole message. This is why the definition talks about the \textbf{message}, not about the key.

  \item \textbf{The message is hard to compute.} Same argument: an adversary that recovers $m$ from $c$ chooses $m_0 \neq m_1$, recovers the message and compares.

  \item \textbf{Every single bit of the message is hard to compute.} Suppose an adversary could compute the first bit of $m$ from $c$ (or the last, or any other bit $i$). Then it chooses $m_0$ and $m_1$ that differ \textbf{exactly in bit $i$}, computes that bit from $c$, and knows which message was encrypted.
\end{itemize}

More generally, the same argument works for \textbf{any} property of the message that can be computed efficiently: choose $m_0$ with the property and $m_1$ without it. Indistinguishability implies that the ciphertext leaks \textbf{no partial information} on the message (except its length). This is the notion of \textit{semantic security}, which is known to be equivalent to indistinguishability (Goldwasser and Micali).

\subsection{Security of the scheme}

\begin{framedthm}[label={thm:prg-otp}]{Encryption from a PRG is one-time secure}
  If $G : \{0,1\}^n \to \{0,1\}^{2n}$ is a PRG, then the scheme of \Cref{def:prg-otp} is one-time computationally secure.
\end{framedthm}

The proof uses \textbf{game hopping}, a very common technique: to compare two games that are far apart, we go from one to the other through a chain of intermediate games, each \textbf{close} to the next one. It is the same idea as the hybrid argument of \Cref{thm:stretch}, with few steps.

\paragraph{The games.} The two games to compare are
\[
  G_0 = \mathrm{GAME}^{\text{one-time}}_{\Pi,\Adv}(\lambda, 0) \qquad\qquad G_1 = \mathrm{GAME}^{\text{one-time}}_{\Pi,\Adv}(\lambda, 1)
\]
and we must show $G_0 \cind G_1$. We introduce two \textbf{hybrid games} $H_0$ and $H_1$: $H_b$ is identical to $G_b$, except that the pad $G(k)$ is replaced by a \textbf{truly uniform} string $z \sample \{0,1\}^{2n}$.

\begin{table}[H]
  \centering
  \small
  \begin{tabular}{@{}lll@{}}
    \toprule
    \textbf{Game} & \textbf{Pad} & \textbf{Ciphertext} \\
    \midrule
    $G_b$ (real game) & $G(k)$ with $k \sample \{0,1\}^n$ & $c = G(k) \oplus m_b$ \\
    $H_b$ (hybrid game) & $z \sample \{0,1\}^{2n}$ & $c = z \oplus m_b$ \\
    \bottomrule
  \end{tabular}
  \caption{In the hybrid games the PRG is replaced by true randomness: $H_b$ is the One-Time Pad.}
\end{table}

\paragraph{The plan.} We go from $G_0$ to $G_1$ in three hops:
\[
  G_0 \;\overset{\text{PRG}}{\cind}\; H_0 \;\overset{\text{OTP}}{\equiv}\; H_1 \;\overset{\text{PRG}}{\cind}\; G_1
\]

\begin{proof}
  \leavevmode
  \textbf{Hop 1: $H_0 \equiv H_1$, by the perfect secrecy of the One-Time Pad.} In $H_b$ the ciphertext is $z \oplus m_b$ with $z$ uniform, that is, a One-Time Pad encryption of $m_b$. By \Cref{prop:otp} it is uniform over $\{0,1\}^{2n}$ whatever the message, so the view of $\Adv$ has the same distribution in $H_0$ and in $H_1$, and $\Pr[H_0 = 1] = \Pr[H_1 = 1]$.

  \medskip
  \textbf{Hop 2: $G_b \cind H_b$, by the security of the PRG.} Assume not: there is a PPT adversary $\Adv$ and a bit $b$ such that $\Adv$ tells $G_b$ and $H_b$ apart with probability at least $\frac{1}{\poly(\lambda)}$. We build a PPT distinguisher $B$ for the PRG. $B$ receives from its challenger a string $z \in \{0,1\}^{2n}$, which is either pseudorandom, $z = G(k)$ with $k \sample \{0,1\}^n$, or truly uniform, $z \sample \{0,1\}^{2n}$. It then plays the role of the challenger for $\Adv$:

  \begin{center}
    \begin{tikzpicture}[x=1cm, y=1cm]
      \node (A) at (0,0) {$\Adv$};
      \node (B) at (4.5,0) {$B$};
      \node (C) at (9,0) {$C$};
      \node[lbl, align=left, anchor=north west] at (9.3,-0.35) {$z = G(k)$, $k \sample \{0,1\}^n$\\ or $z \sample \{0,1\}^{2n}$};
      \draw[msg] (9,-0.6) -- node[above] {$z$} (4.5,-0.6);
      \draw[msg] (4.5,-1.3) -- node[above] {``start''} (0,-1.3);
      \draw[msg] (0,-2.0) -- node[above] {$m_0, m_1$} (4.5,-2.0);
      \draw[msg] (4.5,-2.7) -- node[above] {$c = m_b \oplus z$} (0,-2.7);
      \draw[msg] (0,-3.4) -- node[above] {$b'$} (4.5,-3.4);
      \draw[msg] (4.5,-3.4) -- node[above] {$b'$} (9,-3.4);
    \end{tikzpicture}
  \end{center}

  $B$ runs $\Adv$, receives $m_0, m_1$, answers with $c = m_b \oplus z$, and outputs whatever $\Adv$ outputs. This is a \textbf{perfect simulation}:
  \begin{itemize}
    \item if $z = G(k)$, then $c = G(k) \oplus m_b$ and $\Adv$ is playing \textbf{exactly} $G_b$;
    \item if $z$ is uniform, then $c = z \oplus m_b$ and $\Adv$ is playing \textbf{exactly} $H_b$.
  \end{itemize}
  Therefore
  \[
    \Pr[B \to 1 \mid z = G(k)] = \Pr[\Adv \to 1 \text{ in } G_b] \qquad\qquad \Pr[B \to 1 \mid z \sample \Uni{2n}] = \Pr[\Adv \to 1 \text{ in } H_b]
  \]
  and the advantage of $B$ against the PRG is \textbf{exactly} the advantage of $\Adv$ in telling $G_b$ from $H_b$, at least $\frac{1}{\poly(\lambda)}$. $B$ is PPT, because it runs $\Adv$ and computes one XOR. This contradicts the fact that $G$ is a PRG, so $G_b \cind H_b$ for both values of $b$.

  \medskip
  \textbf{Hop 3: put the hops together.} By the triangle inequality,
  \[
    \begin{aligned}
      \big|\Pr[G_0 = 1] - \Pr[G_1 = 1]\big|
      &\leq \underbrace{\big|\Pr[G_0 = 1] - \Pr[H_0 = 1]\big|}_{\leq\, \negl(\lambda) \text{ (hop 2)}}
      + \underbrace{\big|\Pr[H_0 = 1] - \Pr[H_1 = 1]\big|}_{=\, 0 \text{ (hop 1)}} \\
      &\quad + \underbrace{\big|\Pr[H_1 = 1] - \Pr[G_1 = 1]\big|}_{\leq\, \negl(\lambda) \text{ (hop 2)}} \\
      &\leq 2 \cdot \negl(\lambda)
    \end{aligned}
  \]
  which is negligible. Hence $G_0 \cind G_1$.
\end{proof}

\begin{framedobs}{The structure of the proof}
  \begin{itemize}
    \item \textbf{Each hop changes one thing} $\to$ first the PRG is replaced by true randomness (justified by the PRG assumption), then the message is changed (justified by perfect secrecy), then true randomness is replaced back by the PRG.
    \item \textbf{The reduction is perfect} $\to$ $B$ reproduces the games \textbf{exactly}, so it inherits the whole advantage of $\Adv$, without losses.
    \item \textbf{Computational security in one sentence} $\to$ ``the scheme is as secure as the One-Time Pad, as long as nobody can tell $G(k)$ from a truly random string''.
  \end{itemize}
\end{framedobs}

\paragraph{Looking ahead.} The scheme beats Shannon's bound on the key length, but it is still \textbf{one-time}. Encrypting many messages with the same key requires a stronger definition, in which the adversary sees several ciphertexts, and a scheme that is \textbf{randomized}: a deterministic $\Enc$ always maps the same message to the same ciphertext, so Eve could at least tell when a message is repeated.
